\documentclass[letterpaper,11pt]{article}
\usepackage{graphicx}
\usepackage{geometry}
\linespread{1.2}                        
\setlength{\oddsidemargin}{0.0in}            
\setlength{\textwidth}{6.5in}    
\usepackage[utf8]{inputenc} % Asegúrate de usar utf8
\usepackage[T1]{fontenc} % Asegúrate de usar T1 font encoding      
\setlength{\topmargin}{-.6in}            
\setlength{\textheight}{9in}
\usepackage{latexsym} 
\usepackage{amssymb} 
\usepackage{amsmath}
\usepackage{amsxtra}
\usepackage[mathcal]{eucal} 
\usepackage{enumerate}
\usepackage{hyperref}
\usepackage{rotating}
\usepackage{caption}
\usepackage{lscape}
\usepackage{paralist} 
\usepackage{indentfirst}
\usepackage[dvipsnames,svgnames]{xcolor} 
\usepackage{setspace}
\usepackage{multicol} 
\usepackage{float}

\definecolor{codegreen}{rgb}{0,0.6,0}
\definecolor{codegray}{rgb}{0.5,0.5,0.5}
\definecolor{codepurple}{cmyk}{0.65, 1, 0, 0.35}
\definecolor{backcolour}{rgb}{0.95,0.95,0.92}

\usepackage{listings}
\usepackage{xcolor}

\lstset{
    language=R, % lenguaje de programación
    basicstyle=\small\ttfamily, % estilo básico
    commentstyle=\color{gray}, % estilo para los comentarios
    keywordstyle=\color{blue}, % estilo para las palabras clave
    stringstyle=\color{black}, % estilo para las cadenas de texto
    showstringspaces=false, % no mostrar espacios en cadenas de texto
    tabsize=4, % tamaño de la tabulación
    breaklines=true, % permitir el salto de línea automático
    postbreak=\mbox{\textcolor{red}{$\hookrightarrow$}\space}, % símbolo al final de una línea partida
    numbers=left, % mostrar números de línea
    numberstyle=\tiny\color{gray}, % estilo para los números de línea
    frame=single, % mostrar un borde alrededor del código
    backgroundcolor=\color{gray!10}, % color de fondo
}

\begin{document}

\begin{center}
    {\Large  Taller 1: Econometría }

Profesora:    Jocelyn Tapia Stefanoni\\
Ayudante: Juan Miño
\end{center}

\textbf{Instrucciones:} 

\begin{itemize}
    \item Usted dispone de 40 minutos  para intentar los 100 puntos del taller.
    \item Cuide la presentación y redacción de sus respuestas.
    \item Puede utilizar su computador y los apuntes tanto de clase como de ayudantía.
    \item Debe enviar un archivo en formato PDF y un R script (extensión .R)
\end{itemize}

\section{Datos Aleatorios (40 puntos)}
\begin{enumerate}
    \item Genere el dataframe "datos" utilizando el código en Anexos. (8 pts.)

Según la letra por la que comience su apellido, cambie (o mantenga) la semilla:

\begin{table}[htbp]
\centering
\begin{tabular}{|c|c|c|c|c|}
\hline
\textbf{A-E} & \textbf{F-J} & \textbf{K-O} & \textbf{P-T} & \textbf{U-Z} \\
\hline
123 & 456 & 789 & 101112 & 131415 \\
\hline
\end{tabular}
% \caption{División del alfabeto en 5 segmentos con semillas asignadas}
\label{tabla:division-alfabeto}
\end{table}

El set de datos generado explica las notas de 50 alumnos de la USM:

\begin{itemize}
    \item \textbf{Horas de sueño}: Horas de sueño promedio que cada alumno ha tenido en el último mes.
    \item \textbf{Profesor particular}: Indica si cada alumno ha recibido ayuda externa de un profesor particular o no.
    \item \textbf{Promedio notas semestre pasado}: Promedio de notas que cada alumno obtuvo en el semestre pasado.
    \item \textbf{Tiempo de estudio}: Tiempo de estudio en horas que cada alumno dedica.
    \item \textbf{Asistencia a clases}: Porcentaje de asistencia a clases de cada alumno.
    \item \textbf{Nivel socioeconomico}: Nivel socioeconómico de cada alumno, con valores del 1 al 5.
\end{itemize}

A partir del dataset generado:
    \item Redondee las notas a 2 decimales. (4 pts.)

\begin{lstlisting}
datos$notas <- pmax(pmin(datos$notas, 100), 20)
\end{lstlisting}

    \item Calcule manualmente los beta estimadores a través de la matriz X y el vector de la variable dependiente. (8 pts.)

\begin{lstlisting}
# Calcular X
X <- cbind(rep(1, nrow(datos)), datos$hrs_sueno, datos$profesor_part, datos$media_sem_pasado, datos$tiempo_est, datos$asistencia, datos$nivel_socioec)

# Calcular X^T
Xt <- t(X)

# Calcular X^T X
XtX <- t(X) %*% X

# Calcular X^T y
Xty <- t(X) %*% datos$notas

# Calcular los estimadores beta
beta_hat <- solve(XtX) %*% Xty
\end{lstlisting}

    \item Genere un \textit{summary} a partir de un modelo de regresión múltiple, esta vez utilizando funciones en R. (10 pts.)
\begin{lstlisting}
# Generar modelo utilizando función LM
nombres_vars <- c("hrs_sueno", "profesor_part", "media_sem_pasado", "tiempo_est", "asistencia", "nivel_socioec")
modelo <- lm(datos$notas ~ ., data = datos[, nombres_vars])
summary(modelo)

# Ver el resumen del modelo
summary(modelo)
\end{lstlisting}
    \item Utilice la biblioteca \textit{ggcorrplot} y analice los resultados. ¿Existe alguna relación entre la matriz de correlación y los $\beta$ calculados? (10 pts.)
\begin{lstlisting}
# Generar matriz de correlación y comparar
matriz_cor <- cor(datos)
ggcorrplot(matriz_cor, type = "full", lab = TRUE)
\end{lstlisting}

\end{enumerate}

\section{Wooldridge  (60 puntos)}

Utilice los datos en el archivo WAGE1 de Wooldridge para determinar los estimadores asociados al  siguiente modelo de regresión poblacional.

$wage=\beta_0+ \beta_1educ +\beta_2exper+ \beta_3tenure+u$


Donde wage es salario, educ son los años de escolaridad, 
exper son los años de experiencia laboral, y tenure los años en su actual trabajo.

\begin{enumerate}
\item Determine los estimadores de MCO. Tienen los signos que usted esperaba?
\item ¿Qué estimadores son significativos estadísticamente?
\item Compare el $R^2$ y el $\bar{R}^2$.
\item Determine la matiz de correlación de las variables de su base datos.
\end{enumerate}

\begin{enumerate}
    \item Utilizando R, carge la base de datos "\textit{ceosal1}" desde Wooldridge. (10 pts.)

\begin{lstlisting}
# Wooldridge
library(wooldridge)
data("ceosal1")
ceosal1 <- na.omit(ceosal1)
\end{lstlisting}

    \item Sub seleccione las siguientes variables y almacénelas en un nuevo set de datos (12 pts.):

\begin{itemize}
\item \textbf{pcsalary}: Porcentaje de cambio en el salario del CEO.
\item \textbf{sales}: Ventas totales de la empresa.
\item \textbf{roe}: Retorno sobre el patrimonio (ROE).
\item \textbf{pcroe}: Porcentaje de cambio en el ROE.
\item \textbf{ros}: Retorno sobre las ventas (ROS).
\item \textbf{indus}: Indicador de la industria (1 si es del sector industrial, 0 en otro caso).
\item \textbf{finance}: Indicador del sector financiero (1 si es del sector financiero, 0 en otro caso).
\item \textbf{consprod}: Indicador del sector de bienes de consumo (1 si es del sector de bienes de consumo, 0 en otro caso).
\item \textbf{utility}: Indicador del sector de servicios públicos (1 si es del sector de servicios públicos, 0 en otro caso).
\item \textbf{lsalary}: Logaritmo natural del salario del CEO.
\item \textbf{lsales}: Logaritmo natural de las ventas totales de la empresa.
\end{itemize}

\begin{lstlisting}
columnas_deseadas <- c("pcsalary", "sales", "roe", "pcroe", "ros", "indus", "finance", "consprod", "utility", "lsalary", "lsales")
ceosal1 <- ceosal1[columnas_deseadas]
\end{lstlisting}

    \item Genere un modelo de regresión lineal múltiple que contenga las variables anteriores y un resumen del modelo utilizando \textit{summary}. (12 pts.)
\begin{lstlisting}
# Dividir el conjunto de datos en variables independientes (predictoras) y la variable dependiente (salary)
X <- ceosal1[, c("pcsalary", "sales", "roe", "pcroe", "ros", "indus", "finance", "consprod", "utility", "lsalary", "lsales")]
y <- ceosal1$salary

# Ajustar un modelo de regresion
model <- lm(y ~ ., data = X)

# Resumen del modelo
summary(model)
\end{lstlisting}

    \item Genere un nuevo modelo pero excluyendo las variables no estadísticamente significativas (\textit{p-value < 0.05}). (12 pts.)
\begin{lstlisting}
# Dividir el conjunto de datos en variables independientes (predictoras) y la variable dependiente (salary)
X2 <- ceosal1[, c("sales", "roe", "indus", "finance", "lsalary", "lsales")]
y2 <- ceosal1$salary

# Ajustar un modelo de regresión lineal múltiple
model2 <- lm(y2 ~ ., data = X2)
\end{lstlisting}

    \item Compare su nuevo modelo con su anterior modelo en términos de significancia global y bondad de ajuste. (14 pts.)

\end{enumerate}
\newpage
\section{Anexos: Código}
\begin{verbatim}
# Establecer la semilla para reproducibilidad
set.seed(123)

# Crear la base de datos ficticia con 50 alumnos
datos <- data.frame(
  hrs_sueno = round(runif(50, min = 5, max = 10), 1),
  profesor_part = sample(c(0, 1), 50, replace = TRUE),
  media_sem_pasado = round(runif(50, min = 60, max = 100), 1),
  tiempo_est = round(runif(50, min = 1, max = 8), 1),
  asistencia = round(runif(50, min = 60, max = 100), 1),
  nivel_socioec = sample(1:5, 50, replace = TRUE)
)

# Calcular las notas con ajustes en la ponderación
datos$notas <- 30 + # Base más baja para empezar desde un mínimo más realista
  datos$hrs_sueno * 1.5 + # Reducir el impacto del sueño
  datos$profesor_part * 3 + # Menos puntos adicionales por participación del profesor
  datos$media_sem_pasado * 0.2 + # Menor peso de la media semestral pasada
  datos$tiempo_est * 2 + # Ajuste en tiempo de estudio
  datos$asistencia * 0.15 + # Menor impacto de la asistencia
  datos$nivel_socioec * 2 + # Menos impacto del nivel socioeconómico
  rnorm(50, mean = 0, sd = 5) # Menor desviación estándar para reducir variabilidad

# Asegurar que las notas estén entre 20 y 100
datos$notas <- pmax(pmin(datos$notas, 100), 20)
\end{verbatim}
\end{document}






